%% kapa: exam=1 round=24 year=2013 period=1 form=A subject=civil_law subject_ko=민법 seq=16 number=16
%% source: Q-Net 감정평가사 기출문제 (cst00309) artlSeq=5073802 file=extracted/1008/5073802/2013년 제24회 감정평가사 1차시험 기출문제/1교시(A형).hwp
%% conversion: kordoc 4.24.51 -> tools/gen1.py (auto); review=proofread-sonnet; flags=none
\begin{question}{16}
조건에 관한 설명으로 옳지 않은 것은? (다툼이 있으면 판례에 의함)
\begin{choices}
  \item 조건이 불능인 법률행위는 그 조건이 정지조건이면 무효이고, 해제조건이면 조건 없는 법률행위로 된다.
  \item 해제조건이 법률행위 당시 기성조건이면 그 법률행위는 무효이다.
  \item 조건을 붙일 수 없는 법률행위에 조건을 붙인 경우 그 법률행위는 전부무효임이 원칙이다.
  \item 불법조건이 붙은 법률행위는 정지조건이면 무효이고, 해제조건이면 조건 없는 행위로서 유효하다.
  \item 소유권유보부 매매의 경우 대금이 모두 지급되면 별도의 의사표시가 없더라도 목적물의 소유권이 매수인에게 이전된다.
\end{choices}
\end{question}
\answer{4}
